%%%%%%%%%%%%%%%%%%%%%%%%%%%%%%%%%%%%%%%%%%%%%%%%%%%%%%%%%%%%%%%%%%%%%%
% Lớp 11 | Đường thẳng, mặt phẳng. Quan hệ song song trong | Hai đường thẳng song song
% Dạng 1H4?2-1: Câu hỏi lý thuyết
% Số câu: 27  (sinh tự động bởi tach_cau.py; từ nay thầy sửa trực tiếp file này)
%%%%%%%%%%%%%%%%%%%%%%%%%%%%%%%%%%%%%%%%%%%%%%%%%%%%%%%%%%%%%%%%%%%%%%

%%%%%%%%%%%%%%%% TRẮC NGHIỆM 4 PHƯƠNG ÁN %%%%%%%%%%%%%%%%

% nguồn: Ok_Toan11_TracNghiem_Lan2
\begin{ex}%[1H4N2-1]%[ma:56f9135a]
	Trong  không gian, cho ba đường thẳng $a$, $b$, $c$. Trong các mệnh đề sau mệnh đề nào đúng?
	 \choice
		{Nếu $a$ và $b$ cắt nhau, $b$ và  $c$ cắt nhau thì $a$ và $c$ cắt nhau}
		{Nếu $a$ và  $b$ không cắt nhau thì $a$ và $b$ song song}
		{Nếu $a$ và $b$ cùng chéo nhau với $c$ thì $a$ song song  với $b$}
		{\True Nếu $b$ và $c$ chéo nhau thì $b$ và $c$  không cùng thuộc một mặt phẳng}
	 \loigiai{
		Phương án \lq\lq Nếu $a$ và $b$ cắt nhau, $b$ và $c$ cắt nhau thì $a$ và $c$ cắt nhau \rq\rq~sai vì:  Nếu $a$ và $b$ cắt nhau, $b$ và $c$ cắt nhau thì $a$ và $c$ cắt nhau, chéo nhau hoặc song song.\\
		Phương án \lq\lq Nếu $a$ và $b$ không cắt nhau thì $a$ và $b$ song song \rq\rq~sai vì: Nếu $a$ và $b$ không cắt  nhau thì $a$ và $b$ song song hoặc chéo nhau.\\
	Phương  án \lq\lq Nếu $a$ và $b$ cùng chéo nhau với $c$ thì $a$ song song với $b$ \rq\rq~sai vì: Nếu $a$ và $b$ cùng chéo nhau với $c$ thì $a$ và $b$ song song hoặc cắt nhau.}
\end{ex}

% nguồn: Ok_Toan11_TracNghiem_Lan2
\begin{ex}%[1H4N2-1]%[ma:e701f0cc]
	Trong các mệnh đề  sau, mệnh đề nào đúng?
	 \choice
		{Hai đường thẳng không có  điểm chung thì song song}
		{Hai đường thẳng lần lượt nằm  trên hai mặt phẳng phân biệt thì chéo nhau}
		{Hai đường thẳng phân biệt không song song thì  chéo nhau}
		{\True Hai đường thẳng song  song thì không có điểm chung}
	 \loigiai{
	 Trong 4 mệnh đề trên thì mệnh đề đúng là \lq\lq Hai đường thẳng song song thì không có điểm chung\rq\rq.}
\end{ex}

% nguồn: Ok_Toan11_TracNghiem_Lan2
\begin{ex}%[1H4N2-1]%[ma:b5bd8f71]
	Mệnh đề nào dưới  đây đúng?
	\choice 
		{Trong không gian, hai đường thẳng cùng song song với đường  thẳng thứ ba thì song song với nhau}
		{\True Trong không gian, hai đường thẳng phân biệt cùng song song với đường thẳng thứ ba thì song song với nhau} 
		 {Trong không gian, hai đường thẳng phân biệt cùng song song với đường thẳng thứ ba thì chéo nhau}
		{Trong không gian, hai đường thẳng cùng song song  với đường thẳng thứ ba thì chéo nhau}
	\loigiai{ 
	Theo định lí 3 trang 104 SGK11 Chân trời sáng tạo  thì \lq\lq Hai đường thẳng phân biệt cùng song song với đường thẳng thứ ba thì song song với nhau\rq\rq.}
\end{ex}

% nguồn: Ok_TracNghiemLop11_Lan1
\begin{ex}%[1H4N2-1]%[ma:bed13774]
	Cho các mệnh đề sau đây, mệnh đề  nào đúng?
	\choice 
		{Trong  không gian hai đường thẳng phân biệt không chéo nhau thì cắt nhau}
		{Trong không gian hai đường thẳng phân  biệt không song song thì chéo nhau}
		{Trong  không gian hai đường thẳng phân biệt không có điểm chung thì chéo nhau}
		{\True Trong không gian hai đường thẳng phân biệt cùng nằm trong một mặt  phẳng thì không chéo nhau}
	 \loigiai{
		Mệnh đề đúng là: \lq\lq  Trong không gian hai đường thẳng phân biệt cùng nằm trong một mặt phẳng thì không chéo nhau\rq\rq.
	 }
\end{ex}

% nguồn: Ok_TracNghiemLop11_Lan1
\begin{ex}%[1H4N2-1]%[ma:2b401890]
	Trong không gian, hai mặt phẳng song song với nhau khi và  chỉ khi
	 \choice
		{Có một mặt phẳng chứa hai đường thẳng phân biệt cùng song song với mặt  phẳng còn lại}
		{Hai mặt  phẳng cùng song song với một đường thẳng}
		{Hai mặt phẳng cùng song song  với mặt phẳng thứ ba}
		{\True Hai mặt phẳng  không có điểm chung}
	 \loigiai{
		Trong không gian, hai mặt phẳng song song với nhau khi và chỉ khi hai mặt phẳng không có  điểm chung.
	 }
\end{ex}

% nguồn: Ok_TracNghiemLop11_Lan1
\begin{ex}%[1H4N2-1]%[ma:af9bb978]
	Cho hai đường thẳng chéo nhau  $a$ và $b$. $(P)$ chứa $a$ và song song với $b$, $(Q)$ chứa $b$ và song song với $a$. Phát biểu nào sau đây là đúng?
	 \choice
		{$(P)$  và $(Q)$ cắt nhau}
		{\True $(P)$ và $(Q)$ song song với  nhau}
		 {$(P)$ và $(Q)$ trùng nhau}
		{$(P)$ và $(Q)$ cắt nhau hoặc song song với nhau} 
	\loigiai{ 
		Cho hai  đường thẳng chéo nhau $a$ và $b$. $(P)$ chứa $a$ và song song với $b$, $(Q)$ chứa $b$ và song song với $a$, khi đó $(P)$ và $(Q)$ song song với nhau.
	 }
\end{ex}

% nguồn: Ok_TracNghiemLop11_Lan1
\begin{ex}%[1H4N2-1]%[ma:d0e89cd0]
	Cho hai đường thẳng phân biệt không có điểm chung cùng nằm trong một mặt phẳng  thì hai đường thẳng đó
	\choice 
		{trùng  nhau}
		{chéo nhau} 
		{\True  song song}
		{cắt  nhau}
	 \loigiai{
	Hai đường thẳng phân biệt không có điểm chung cùng nằm trong một mặt phẳng  thì hai đường thẳng đó song song.}
\end{ex}

% nguồn: Ok_TracNghiemLop11_Lan1
\begin{ex}%[1H4N2-1]%[ma:19a418aa]
	Cho $d\parallel (\alpha)$, mặt phẳng $(\beta)$ qua $d$ cắt  $(\alpha)$ theo giao tuyến $d'$. Khẳng định nào sau đây là đúng?
	 \choice
		 {$d$ cắt $d'$}
		{\True  $d\parallel d' $}
		{$d$ và  $d'$ chéo nhau}
		{$d\equiv d' $} 
	 \loigiai{
		Ta có $d'=(\alpha)\cap (\beta)$. Do $d$ và $d'$ cùng thuộc $(\beta)$ nên  $d$ cắt $d'$ hoặc $d\parallel d'$.\\
		Nếu $d$  cắt $d'$ thì $d$ cắt $(\alpha)$ (mâu thuẫn với giả thiết).
	Vậy $d\parallel d'$.} 
\end{ex}

% nguồn: Ok_TracNghiemLop11_Lan1
\begin{ex}%[1H4N2-1]%[ma:df2a5011]
	Chọn khẳng định đúng trong các khẳng định sau. 
	 \choice
		{Trong không gian, hai đường thẳng cắt nhau là  hai đường thẳng không có điểm chung}
		{Trong không gian,  hai đường thẳng không có điểm chung là hai đường thẳng song song}
		{Trong không gian, hai đường thẳng  không có điểm chung là hai đường thẳng chéo nhau}
		{\True Trong không gian, hai đường  thẳng song song là hai đường thẳng cùng nằm trong một mặt phẳng và không có điểm chung}
	\loigiai{ Trong không gian, 
		\begin{enumerate}[\faCheckSquareO]
			\item Hai đường thẳng cắt nhau là hai đường thẳng có một điểm chung duy nhất. 
			\item Hai  đường thẳng song song là hai đường thẳng cùng nằm trong một mặt phẳng và không có điểm chung
			\item Hai đường thẳng chéo nhau là  hai đường thẳng không cùng nằm trên một mặt phẳng. Suy ra hai đường thẳng chéo nhau không có điểm chung.
		\end{enumerate}
		Vậy trong không gian, hai  đường thẳng không có điểm chung có thể là hai đường thẳng song song hoặc hai đường thẳng chéo nhau.
	} 
\end{ex}

% nguồn: Ok_TracNghiemLop11_Lan1
\begin{ex}%[1H4N2-1]%[ma:ec7d4509]
	Chọn mệnh đề đúng trong các mệnh đề sau. 
	\choice 
		{Hai đường thẳng không có điểm chung thì  chéo nhau}
		{Hai đường thẳng  phân biệt không có điểm chung thì chéo nhau}
		{\True Hai  đường thẳng chéo nhau thì không có điểm chung}
		{Hai đường thẳng lần  lượt nằm trên hai mặt phẳng phân biệt thì chéo nhau}
	 \loigiai{
		- Hai đường  thẳng chéo nhau thì không có điểm chung là \textbf{đúng}.
		- Hai đường thẳng không có điểm chung thì chéo nhau \textbf{hoặc song  song nhau}.
		 - Hai đường thẳng lần lượt nằm trên hai mặt phẳng phân biệt thì chéo nhau hoặc song song nhau cũng có thể cắt nhau.
	 }
\end{ex}

% nguồn: Ok_TracNghiemLop11_Lan1
\begin{ex}%[1H4N2-1]%[ma:0d894af6]
	Chọn khẳng  định {\bf sai}.
	 \choice
		{Hai đường thẳng  phân biệt cùng song song với đường thẳng thứ ba thì song song với nhau}
		 {Nếu hai đường thẳng chéo nhau thì chúng không đồng phẳng}
		{\True Hai đường thẳng song song thì  không đồng phẳng và không có điểm chung}
		{Hai đường thẳng cắt nhau thì đồng phẳng và có một điểm chung} 
	\loigiai{ 
		Khẳng định sai là \lq\lq Hai đường  thẳng song song thì không đồng phẳng và không có điểm chung\rq\rq.
	} 
\end{ex}

% nguồn: Ok_TracNghiemLop11_Lan1
\begin{ex}%[1H4N2-1]%[ma:b163b7d0]
	Trong các mệnh đề sau, mệnh  đề nào \textbf{sai}?
	 \choice
		{\True Hai đường thẳng không có  điểm chung thì chéo nhau}
		{Hai đường thẳng chéo  nhau thì không có điểm chung}
		{Hai đường thẳng phân biệt không cắt nhau và không song song thì chéo nhau} 
		{Hai đường thẳng phân biệt không  chéo nhau thì hoặc cắt nhau hoặc song song}
	\loigiai{ 
		Hai đường thẳng không có điểm chung thì  chúng song song (khi chúng đồng phẳng) hoặc chéo nhau (khi chúng không đồng phẳng).
	 }
\end{ex}

% nguồn: Ok_Toan11_TracNghiem_Lan2
\begin{ex}%[1H4H2-1]%[ma:158fee56]
	Cho hai đường thẳng chéo nhau $a$ và $b$. Lấy điểm $A$, $B$ thuộc $a$ và $C$, $D$ thuộc $b$. Khẳng định nào sau đây đúng khi nói về hai  đường thẳng $AD$ và $BC$?
	\choice 
		{Có thể song song hoặc cắt nhau} 
		{Cắt  nhau}
		{Song song  nhau}
		{\True Chéo  nhau}
	 \loigiai{
	\immini{Theo giả thiết, $a$ và $b$ chéo nhau $\Rightarrow a$ và $b$ không đồng phẳng.\\ 
		Giả sử $AD$ và $BC$  đồng phẳng.\\
		Nếu $AD\cap BC=I\Rightarrow I\in (ABCD)\Rightarrow I\in mp(a;b)$.  Mà $a$ và $b$ không đồng phẳng, do đó, không tồn tại điểm $I$.\\
		Nếu $AD\parallel BC\Rightarrow  a$ và $b$ đồng phẳng.\\
		Vậy điều giả sử là sai. Do đó $AD$ và $BC$ chéo nhau.}{\begin{tikzpicture}[
			 x=1cm,y=1cm,
			 xscale=1,yscale=-1,
			 solid/.style={draw,line cap=round,very thick},
			dashed/.style={solid,dash  pattern=on 0.2cm off 0.2cm},
			point/.style={solid,fill,circle,inner sep=1pt} 
			] 
		\coordinate  (A) at (5.9,2.9);
		\coordinate (B) at  (8.2,1.5);
		\draw[solid]  (4.1,1.1) -- (3.3,4.1) -- (8.0,3.4) -- (9.1,3.2);
		\draw[solid] (4.6,4.7)  -- (3.8,7.1) -- (9.1,6.2);
		\draw[solid]  (8.0,3.4) -- (7.9,5.6);
		\draw[dashed] (8.2,1.5)  -- (8.0,3.4);
		\draw[solid] (6.1,3.7) --  (6.7,5.4);
		\draw[dashed]  (5.9,2.9) -- (6.1,3.7);
		 \draw[very thick]
			($(A)!-0.4!(B)$)  -- ($(B)!-0.3!(A)$);
		\node[point,label={above left:$A$}] at (A)  {};
		\node[point,label={above:$B$}] at (B)  {};
		\coordinate (C)  at (7.9,5.6);
		\coordinate (D) at  (6.7,5.4);
		 \draw[very thick]
			($(C)!-0.5!(D)$) --  ($(D)!-1!(C)$);
		 \node[point,label={below:$C$}] at (7.9,5.6) {};
		 \node[point,label={below:$D$}] at (6.7,5.4) {};
		\node[label={left:$a$}] at (5.3,3.2)  {};
		\node[label={above:$b$}] at (5.0,5.27)  {};
	\end{tikzpicture}
		 }
	 }
\end{ex}

% nguồn: Ok_TracNghiemLop11_Lan1
\begin{ex}%[1H4H2-1]%[ma:14ecf502]
	Cho hình chóp $S.ABCD$ có đáy $ABCD$ là hình bình hành tâm $O$. Gọi $M$, $N$ lần lượt là trung điểm của $SA$  và $SD$. Trong các khẳng định sau, khẳng định nào \textbf{sai}?
	\choice 
		{$MN \parallel  BC$}
		 {$ON \parallel SB$}
		{$OM  \parallel SC$}
		{\True$ON  \parallel SC$}
	 \loigiai{
		\begin{center}
			\begin{tikzpicture}[line join = round,line cap = round,font=\footnotesize,>=stealth,scale=1]
					\coordinate (B) at  (0,0);
					 \coordinate (C) at (5,0);
					\coordinate (A) at  (2,2);
					 \coordinate (D) at ($(A)-(B)+(C)$);
					\coordinate (S) at  (1,5);
					\coordinate  (M) at ($(S)!0.5!(A)$);
					\coordinate  (N) at ($(S)!0.5!(D)$);
					\coordinate (O)  at ($(A)!0.5!(C)$);
					\draw  (B)--(C)--(D)--(S)--(B)--(S)--(C);
					\draw[dashed]  (D)--(B)--(A)--(D)
						(S)--(A)--(C) 
						(M)--(N)--(O)--(M); 
					\foreach \i/\g in {A/90,B/-90,C/-90,D/90,S/90,M/180,N/90,O/-90}{\draw[fill](\i) circle (1pt) 
						 ($(\i)+(\g:3mm)$) node[scale=1]{$\i$};}
				\end{tikzpicture}
		\end{center}
		Ta có $M N$ là đường trung bình của tam giác $S A D$. Suy ra, $M N \parallel A D$  mà $A D \parallel B C$ nên $M N \parallel B C$.\\
		Xét tam giác $S B D$ có  $O N$ là đường trung bình. Suy ra, $O N \parallel S B$.\\
		Xét tam giác $S A C$ có $O M$ là đường  trung bình. Suy ra, $O M \parallel S C$.\\
		Vậy khẳng định $O N \parallel  S C$ là sai.
	 }
\end{ex}

% nguồn: Ok_TracNghiemLop11_Lan1
\begin{ex}%[1H4H2-1]%[ma:5cb6e277]
	Cho tứ diện $ABCD$. Khi đó hai đường thẳng $AB$  và $CD$ là hai đường thẳng
	\choice 
		{cắt  nhau}
		 {song song}
		 {\True chéo nhau}
		 {trùng nhau}
	 \loigiai{
		Vì  $ABCD$ là tứ diện nên $AB$ và $CD$ không đồng phẳng.
	} 
\end{ex}

% nguồn: Ok_TracNghiemLop11_Lan1
\begin{ex}%[1H4H2-1]%[ma:ff1f4bf4]
	Khẳng định nào sau đây là \textbf{sai}? 
	 \choice
		{\True Qua ba điểm phân biệt  xác định duy nhất một mặt phẳng}
		{Qua một đường thẳng và một điểm không thuộc đường thẳng xác định duy nhất  một mặt phẳng}
		{Qua hai đường  thẳng song song xác định duy nhất một mặt phẳng}
		{Qua hai đường thẳng cắt nhau xác định duy nhất một mặt phẳng} 
	 \loigiai{
		Khẳng định sai là: \lq\lq Qua ba điểm phân biệt xác định duy nhất một mặt phẳng\rq\rq. Vì ba điểm thẳng  hàng thì có vô số mặt phẳng chứa ba điểm đó.
	 }
\end{ex}


%%%%%%%%%%%%%%%% ĐÚNG/SAI %%%%%%%%%%%%%%%%

% nguồn: Ok_dungsai_lop11_Lan2
\begin{ex}%[1H4N2-1]%[ma:7b35f6d4]
	Cho biết tính đúng sai của mỗi phát biểu sau (xét  trong không gian)
	\choiceTF 
		{Hai đường  thẳng không có điểm chung thì chúng song song với nhau}
		{Hai đường thẳng không có điểm  chung thì chúng chéo nhau}
		{Hai đường thẳng có điểm chung  thì chúng cắt nhau}
		{\True Hai đường thẳng không thể cùng nằm trên  một mặt phẳng thì chúng chéo nhau}
	\loigiai{ 
		\begin{itemchoice}
			\itemch {\bf  Sai}.\\
			Hai đường thẳng  không có điểm chung thì chúng song song với nhau hoặc chéo nhau.
			 \itemch {\bf Sai}.\\
			Hai đường thẳng không có điểm chung thì chúng  chéo nhau hoặc song song với nhau.
			\itemch  {\bf Sai}.\\
			Hai đường thẳng có điểm chung thì chúng  cắt nhau hoặc trùng nhau.
			\itemch {\bf  Đúng}.\\
			Hai đường  thẳng không thể cùng nằm trên một mặt phẳng thì chúng chéo nhau.
		\end{itemchoice}
	 }
\end{ex}

% nguồn: Ok_dungsai_lop11_Lan2
\begin{ex}%[1H4N2-1]%[ma:baff5f76]
	Trong không gian cho ba đường  thẳng $a$, $b$ và $c$ phân biệt. Trong các phát biểu sau, phát biểu nào đúng, phát biểu nào \textbf{sai}?
	\choiceTF 
		{Nếu hai đường thẳng cùng song  song với đường thẳng thứ ba thì chúng song song với nhau}
		{Nếu hai đường  thẳng cùng chéo nhau với đường thẳng thứ ba thì chúng chéo nhau}
		{\True Nếu đường thẳng $a$ song song với đường thẳng  $b$, đường thẳng $b$ và đường thẳng $c$ chéo nhau thì đường thẳng $a$ và đường thẳng $c$ chéo nhau hoặc cắt nhau}
		{Nếu đường thẳng $a$ cắt $b$, hai đường thẳng $b$ và  $c$ chéo nhau thì $a$ và $c$ chéo nhau hoặc song song với nhau}
	 \loigiai{
		\begin{itemchoice}
			\itemch {\bf Sai}.\\ 
			Nếu  hai đường thẳng cùng song song với đường thẳng thứ ba thì chúng song song với nhau hoặc trùng nhau.
			\itemch {\bf Sai}.\\ 
			Nếu hai  đường thẳng cùng chéo nhau với đường thẳng thứ ba thì chúng có thể chéo nhau, song song, cắt nhau hoặc trùng nhau.
			 \itemch {\bf Đúng}.\\
			Nếu đường thẳng $a$ song song với đường thẳng  $b$, đường thẳng $b$ và đường thẳng $c$ chéo nhau thì đường thẳng $a$ và đường thẳng $c$ chéo nhau hoặc cắt nhau.
			\itemch  {\bf Sai}.\\
			Đường thẳng $a$ có thể cắt cả hai đường chéo nhau là $b$ và $c$, tức  là đường thẳng $a$ có thể cắt đường thẳng $c$.
		\end{itemchoice}
	 }
\end{ex}

% nguồn: Ok_dungsai_lop11_Lan2
\begin{ex}%[1H4N2-1]%[ma:1773a496]
	Trong không gian, cho ba đường thẳng $a$, $b$, $c$ phân biệt. Xét tính đúng sai  của các mệnh đề sau
	 \choiceTF
		{Nếu $a$ và $b$ không cắt nhau thì  $a$ và $b$ song song}
		{\True Nếu $a$ và $b$ cùng song song với $c$  thì $a$ và $b$ song song}
		{Nếu $a$ và $b$ cắt nhau, $c$ và  $b$ cắt nhau thì $a$ và $c$ cắt nhau}
		{Nếu $a$ và $b$ chéo nhau, $c$ và $b$ chéo nhau thì $a$ và $c$ song song} 
	 \loigiai{
		\begin{itemchoice}
			\itemch Sai. Nếu $a$  và $b$ không cắt nhau thì $a$ và $b$ song song hoặc chéo nhau.
			\itemch Đúng. Nếu $a$ và $b$ cùng song song với $c$ thì $a$ và $b$ song song. 
			\itemch Sai. Nếu $a$ và $b$ cắt nhau, $c$ và $b$ cắt nhau thì $a$ và $c$ cắt nhau, song song hoặc chéo  nhau.
			\itemch Sai. Nếu  $a$ và $b$ chéo nhau, $c$ và $b$ chéo nhau thì $a$ và $c$ song song, chéo nhau hoặc cắt nhau.
		\end{itemchoice}
	} 
\end{ex}

% nguồn: Ok_dungsai_lop11_Lan2
\begin{ex}%[1H4N2-1]%[ma:c8e86aa1]
	Cho hình hộp $ABCD.A'B'C'D'$. 
	 \choiceTF
		{\True $BD'$ và $B'C'$ chéo nhau} 
		{\True $A'C$ và $DD'$ chéo nhau} 
		{$DC'$  và $AB'$ chéo nhau}
		{\True $AB'C'D$ và $A'BCD'$ là hai hình bình hành  có chung một đường trung bình}
	\loigiai{ 
		\begin{center}
			\begin{tikzpicture}[line join=round,line cap=round,>=stealth,font=\footnotesize,scale=1]
					 \def\ab{2.5}
					 \def\ad{4}
					 \def\h{4}
					 \def\doxien{1}
					\path 
						(0,0) coordinate (A) 
						(30:\ab) coordinate (B) 
						 (0:\ad) coordinate (D)
						($(B)+(D)-(A)$) coordinate (C); 
					\foreach \a in {A,B,C,D}{\coordinate (\a')  at ($(\a)+(\doxien,\h)$);}
					\path 
						($(B')!0.5!(A)$) coordinate  (M)
						($(C')!0.5!(D)$)  coordinate (N);
					\draw (A')--(B')--(C')--(D')--(A')--(A)--(D)--(C)--(C') 
						 (D')--(D)
						(C')--(D) 
						 (D')--(C);
					\draw[dashed]  (A)--(B)--(C)
						 (B)--(B')
						 (M)--(N)
						(A')--(B) 
						(A)--(B'); 
					\foreach \i/\g  in {A'/90,B'/90,C'/90,D'/90,A/-90,B/180,C/-90,D/-90,M/180,N/0}{\draw[fill=black](\i) circle (1pt)
						($(\i)+(\g:3mm)$) node[scale=1]{$\i$};} 
				\end{tikzpicture}
		\end{center}
		\begin{itemchoice}
			\itemch  \textbf{Đúng}.
			\itemch  \textbf{Đúng}.
			\itemch \textbf{Sai}.\\ 
			Hai đường  thẳng đó song song với nhau.
			\itemch \textbf{Đúng}.\\ 
			Đường trung bình là đường thẳng $MN$. 
		\end{itemchoice}
	} 
\end{ex}

% nguồn: Ok_dungsai_lop11_Lan2
\begin{ex}%[1H4H2-1]%[ma:4b0b0c6c]
	 Trong không gian cho ba đường thẳng $a$, $ b$ và $c$ phân biệt.
	\choiceTF 
		{\True Nếu hai đường thẳng cùng song song với đường thẳng thứ ba thì chúng  song song với nhau}
		{Nếu hai đường thẳng cùng chéo nhau với đường thẳng  thứ ba thì chúng chéo nhau}
		{\True Nếu đường thẳng $a$ song song với đường thẳng $b$,  đường thẳng $b$ và đường thẳng $c$ chéo nhau thì đường thẳng $a$ và đường thẳng $c$ chéo nhau hoặc cắt nhau}
		{Nếu đường thẳng $a$ cắt $b$, hai đường thẳng $b$ và $c$ chéo  nhau thì $a$ và $c$ chéo nhau hoặc song song với nhau}
	\loigiai{ 
		\begin{itemchoice}
			\itemch Nếu hai đường thẳng cùng song song với đường thẳng thứ ba thì chúng song song với nhau. 
			\itemch Vì nếu hai đường $a$, $c$ chéo nhau và hai đường $b$, $c$ chéo nhau  thì đường thẳng $a$ và đường thẳng $b$ có đến ba khả năng: chéo nhau, song song hoặc cắt nhau.
			\itemch Nếu đường thẳng $a$ song song với đường thẳng $b$, đường thẳng $b$ và đường thẳng $c$ chéo nhau thì đường thẳng  $a$ và đường thẳng $c$ chéo nhau hoặc cắt nhau.
			\itemch Vì đường thẳng  $a$ có thể cắt cả hai đường chéo nhau là $b$ và $c$, tức là đường thẳng $a$ có thể cắt đường thẳng $c$.
		\end{itemchoice}
	} 
\end{ex}

% nguồn: Ok_dungsai_lop11_Lan2
\begin{ex}%[1H4H2-1]%[ma:67d1bf8e]
	Trong không gian cho ba đường thẳng $a$, $b$ và $c$  phân biệt.
	\choiceTF 
		{\True Nếu $a$ và $b$ cùng song song với $c$ thì chúng  song song với nhau}
		{Nếu $a$ và $c$ chéo nhau, đồng thời $b$ và $c$ chéo nhau thì $a$ và $b$ chéo nhau} 
		{\True Nếu $a$ và $b$ song song với nhau, đồng thời $b$ và $c$ chéo nhau thì $a$ và $c$ chéo nhau  hoặc cắt nhau}
		{Nếu $a$ và $b$ cắt nhau, đồng thời $b$ và $c$ chéo nhau thì $a$ và  $c$ chéo nhau hoặc song song với nhau}
	 \loigiai{
		\begin{itemchoice}
			\itemch \textbf{Đúng} \\ 
			Vì đây là định lý. 
			\itemch  \textbf{Sai}\\
			Vì nếu hai đường $a, c$ chéo nhau và hai đường $b,  c$ chéo nhau thì đường thẳng $a$ và đường thẳng $b$ có đến ba khả năng là chéo nhau, song song hoặc cắt nhau.
			\itemch \textbf{Đúng}\\ 
			Dựa theo vị trí tương đối  của hai đường thẳng trong không gian.
			\itemch  \textbf{Sai} \\
			Vì đường thẳng $a$ có thể cắt cả hai đường chéo nhau  là $b$ và $c$.
		\end{itemchoice}
	 }
\end{ex}

% nguồn: Ok_dungsai_lop11_Lan2
\begin{ex}%[1H4H2-1]%[ma:2eec4941]
	Cho biết tính đúng sai của mỗi phát biểu sau (xét trong không gian). 
	 \choiceTF
		{Hai đường thẳng không có điểm chung thì chúng  song song với nhau}
		 {Hai đường thẳng không có điểm chung thì chúng chéo nhau}
		 {Hai đường thẳng có điểm chung thì chúng cắt nhau}
		{\True Hai đường thẳng không thể cùng nằm trên một  mặt phẳng thì chúng chéo nhau}
	\loigiai{ 
		\begin{itemchoice}
			\itemch  Hai đường thẳng không có điểm chung thì chúng song song hoặc chéo nhau.
			\itemch Hai đường thẳng  không có điểm chung thì có thể là chúng chéo nhau hoặc song song với nhau.
			\itemch Hai đường có điểm chung thì chúng có  thể cắt nhau hoặc trùng nhau.
			\itemch Hai đường thẳng không thể cùng nằm trên một mặt phẳng thì chúng chéo nhau theo  tính chất cơ bản.
		\end{itemchoice}
	 }
\end{ex}


%%%%%%%%%%%%%%%% TRẢ LỜI NGẮN %%%%%%%%%%%%%%%%

% nguồn: Ok_TraLoiNgan_Lop11_Lan2
\begin{ex}%[1H4N2-1]%[ma:3a002b95]
	Cho hai đường thẳng phân biệt trong không gian, có thể xảy ra một trong các trường hợp sau\\ 
	(I)  Cắt nhau.\\
	(II) Song song với  nhau.\\
	 (III) Trùng nhau.\\
	(IV) Chéo  nhau.\\
	Có bao nhiêu  khẳng định đúng.
	\shortans{3} 
	\loigiai{ 
	Theo định nghĩa, nếu hai đường thẳng phân biệt trong không gian đồng phẳng thì sẽ có $2$ vị trí tương đối là cắt nhau, song song nhau. Nếu hai đường thẳng phân biệt, không đồng phẳng trong không gian thì chúng  chéo nhau.}
\end{ex}

% nguồn: Ok_TraLoiNgan_Lop11_Lan2
\begin{ex}%[1H4N2-1]%[ma:9ec188f0]
	Cho hai đường thẳng chéo nhau $ a$, $ b$ và điểm $ M$ không thuộc $ a$ cũng  không thuộc $ b$. Có nhiều nhất bao nhiêu đường thẳng đi qua $ M$ và đồng thời cắt cả $ a$ và $ b$?
	\shortans{1} 
	 \loigiai{
		Gọi $(P)$ là mặt phẳng qua $ M$ và chứa $ a$; $(Q)$ là  mặt phẳng qua $ M$ và chứa $ b$.\\
		Giả sử tồn tại đường thẳng $ c$ đi qua $  M$ và đồng thời cắt cả $ a$ và $ b$ suy ra
		$\heva{  &c\in(P)\\
		 &c\in(Q)}
		 \Rightarrow c=(P)\cap(Q)$.\\
		Mặt khác nếu có một đường thẳng $ c'$ đi qua $ M$ và đồng thời cắt cả  $ a$ và $ b$thì $ a$ và $ b$ đồng phẳng.\\
	Do đó có duy nhất một đường thẳng đi qua $ M$ và  đồng thời cắt cả $ a$ và $ b$.}
\end{ex}

% nguồn: Ok_TraLoiNgan_Lop11_Lan1
\begin{ex}%[1H4H2-1]%[ma:e7aa0a67]
	Cho tứ diện $ABCD$. Điểm $I$ và $J$ theo thứ tự là trung điểm của $AD$ và $AC$, $G$ là trọng tâm tam giác $BCD$. Giao tuyến của hai mặt phẳng $(GIJ)$ và $(BCD)$ cắt $BD$ tại $E$, cắt $BC$ tại $F$. Tính tỉ  số $\dfrac{IJ}{EF}$?
	 \shortans{1,5}
	\loigiai{ 
	\immini{$IJ$ là đường trung bình của $\triangle ACD$  nên $IJ \parallel CD$ và $IJ = \dfrac{1}{2} CD.\qquad(1)$\\
		Ta có $\heva{&IJ \parallel CD \\  &IJ \subset(GIJ) \\ &CD \subset(BCD) \\ &G \in(IJG) \cap(BCD)} \Rightarrow \heva{&G x=(GIJ) \cap (BCD) \\ &Gx \parallel IJ \parallel CD}.$
		Suy ra giao tuyến $Gx \equiv EF$. 
		Tứ  giác $IJFE$ có $IJ \parallel FE$ nên là hình thang.
		Vì $G$ là trọng tâm $\triangle  BCD$ và $EF \parallel CD$.\\
		Nên $\dfrac{EF}{CD} = \dfrac{BE}{BD} =  \dfrac{BG}{BK} = \dfrac{1}{3}. \qquad(2)$\\
		Từ $(1)$ và $(2)$, suy ra $\dfrac{IJ}{EF} = \dfrac{\dfrac{1}{2}}{\dfrac{2}{3}} = \dfrac{3}{2} = 1{,}5$.}{	\begin{tikzpicture}[scale=0.7,font=\footnotesize,line join=round,line cap=round,>=stealth]
		 \path
			(0,0) coordinate (B) 
			++(-45:3)  coordinate (C)
			 ($(B)+(5,0)$) coordinate (D)
			($(B)+(1,4)$) coordinate  (A)
			 ($(A)!0.5!(C)$) coordinate (J)
			($(A)!0.5!(D)$)  coordinate (I)
			($(B)!2/3!(C)$)  coordinate (F)
			($(B)!2/3!(D)$)  coordinate (E)
			($(E)!0.5!(F)$) coordinate (G) 
			($(C)!0.5!(D)$)  coordinate (K)
			($(E)!1.5!(F)$) coordinate (X) node[above]  {$x$};
		\foreach  \i in {B,C,D}{\draw (A)--(\i);}
		 \draw
			(B)--(C)--(D) 
			 (X)--(F)--(J)--(I);
		\draw[dashed,thin] 
			 (B)--(D)
			(F)--(E)--(I) 
			(J)--(G)--(I) 
			 (B)--(K);
		\foreach \i/\g  in {A/90,B/180,C/-90,D/0,I/30,J/-150,E/45,F/-90,G/-90,K/-30}
		\fill[black] (\i) circle  (1pt)+(\g:4mm) node[scale=1]{$\i$};
	\end{tikzpicture}
		 }
	 }
\end{ex}

% nguồn: Ok_TraLoiNgan_Lop11_Lan1
\begin{ex}%[1H4H2-1]%[ma:1406ade0]
	Cho tứ diện $ABCD$.  Điểm $I$ và $J$ theo thứ tự là trung điểm của $AD$ và $AC$, $G$ là trọng tâm tam giác $BCD$. Giao tuyến của hai mặt phẳng $(GIJ)$ và $(BCD)$ cắt $BD$ tại $E$, cắt $BC$ tại $F$. Tính tỉ số $\dfrac{EF}{IJ}$ (kết quả làm tròn đến hàng phần trăm).
	\shortans{0,67} 
	 \loigiai{
		\begin{center}
			\begin{tikzpicture}[scale=0.7,font=\footnotesize,line join=round,line cap=round,>=stealth]
					 \path
						(0,0) coordinate (B) 
						 ++(-45:3) coordinate (C)
						($(B)+(5,0)$) coordinate  (D)
						($(B)+(1,4)$) coordinate (A) 
						 ($(A)!0.5!(C)$) coordinate (J)
						($(A)!0.5!(D)$) coordinate  (I)
						($(B)!2/3!(C)$) coordinate (F) 
						 ($(B)!2/3!(D)$) coordinate (E)
						($(E)!0.5!(F)$)  coordinate (G)
						($(C)!0.5!(D)$) coordinate  (K)
						($(E)!1.5!(F)$) coordinate (X)  node[above] {$x$};
					\foreach \i in  {B,C,D}{\draw (A)--(\i);}
					 \draw
						(B)--(C)--(D) 
						 (X)--(F)--(J)--(I);
					 \draw[dashed,thin]
						 (B)--(D)
						(F)--(E)--(I) 
						 (J)--(G)--(I)
						 (B)--(K);
					\foreach  \i/\g in {A/90,B/180,C/-90,D/0,I/30,J/-150,E/45,F/-90,G/-90,K/-30}
					\fill[black] (\i) circle (1pt)+(\g:4mm) node[scale=1]{$\i$}; 
				\end{tikzpicture}
		\end{center}
		$IJ$ là đường trung bình của $\triangle ACD$ nên $IJ \parallel CD$ và $IJ =  \dfrac{1}{2} CD.\quad (1)$\\
		Ta có $\heva{&IJ \parallel CD \\ &IJ \subset(GIJ) \\ &CD \subset(BCD) \\ &G \in(IJG) \cap(BCD)} \Rightarrow \heva{&G  x=(GIJ) \cap (BCD) \\ &Gx \parallel IJ \parallel CD.}$\\
		Suy ra giao tuyến $Gx  \equiv EF$.\\
		Tứ giác $IJFE$ có $IJ \parallel FE$ nên là hình thang.\\ 
		 Vì $G$ là trọng tâm $\triangle BCD$ và $EF \parallel CD$.\\
		Nên $\dfrac{EF}{CD} = \dfrac{BE}{BD} = \dfrac{BG}{BK} =  \dfrac{1}{3}. \quad (2)$\\
		Từ $(1)$ và $(2)$, suy  ra $\dfrac{EF}{IJ} = \dfrac{\dfrac{2}{3}}{\dfrac{1}{2}} = \dfrac{2}{3}\approx 0{,}67$.
	} 
\end{ex}
